\section{Research target, conventions, and logical scope}
\label{hqj:sec:scope}

The preceding report \emph{Quadratic-Spectrum Stieltjes Identities}
\cite{predecessor} proves a first-derivative factorization formula for
polynomially weighted quadratic spectra. It classifies the weights whose
local correction vanishes identically in the mass, and obtains finite
polylogarithmic formulas at all nonpositive integer spectral points.
Its final research agenda asks for higher spectral derivatives,
higher-jet factorization defects, noninteger-weight resonances, and the
second derivative of a product with distinct masses. Those questions,
rather than another first-derivative rederivation, are the starting point
of this article.

The canonical manuscript already contains a large system of
Hurwitz--Stieltjes and Gamma identities \cite{repo}. Its source inventory
and incoming collection are too extensive for a claim that every
preserved result was audited here. The contribution below is an explicit,
proved continuation of the inspected predecessor. Some ingredients are
classical: the Hurwitz Laurent expansion, Gamma and polygamma identities,
Poisson summation for a quadratic spectrum, half-integer Bessel formulas,
and the first-derivative multiplicative-anomaly mechanism
\cite{hurwitz,gamma,polygamma,bessel,order,dowker}. We do not claim
historical priority for those ingredients or for every specialization of
the resulting identities.

\subsection{Four questions and their answers}
The predecessor's question on higher-jet defects is answered by
Theorems~\ref{hqj:thm:local} and~\ref{hqj:thm:compiler}, with a complete
vanishing classification in Theorem~\ref{hqj:thm:rigidity}. Its question
on noninteger polynomial degree receives a joint meromorphic answer and
an exact finite-part conversion in Theorem~\ref{hqj:thm:rho}.

The higher-negative-integer question is answered constructively at the
\emph{second} derivative by Theorem~\ref{hqj:thm:second}. This specifies
an adequate explicit function algebra, but does not show it is the
smallest possible one. The two-mass second-derivative and collision
question is answered for every polynomial weight in
Theorems~\ref{hqj:thm:product}--\ref{hqj:thm:collision}. Higher spectral
orders have a complete factorization-defect calculus and mass-transport
law here, not a claim of a finite Fourier reduction at every order.

\subsection{Notation and branches}
Throughout the one-sided statements, $a,b>0$ and $x_n=n+a$, $n\ge0$.
A polynomial $W(x)=\sum_{d=0}^{D}w_dx^d$ may have complex coefficients.
For bilateral statements, $0<a<1$ and the lattice is $x=n+a$, $n\in\Z$.
The positive squared masses are denoted $u=b^2$ and $v=c^2$.

All logarithms of positive numbers are real. On the right half-plane
we use the principal analytic logarithm, so for $x>0$
\[
\log(x^2+b^2)=\log(x+ib)+\log(x-ib).
\]
A logarithm of a quotient in a complex-shift identity means the
\emph{difference} of the specified logarithms; we do not silently change
branches by combining them. Parameter continuation is always from these
domains. The exponential arguments of polylogarithms in the massive
formulas lie strictly inside the unit disk.

We write $\zeta^{(r)}(s,a)=\partial_s^r\zeta(s,a)$ for order derivatives,
not derivatives in $a$, and use the standard Stieltjes convention
\begin{equation}\label{hqj:eq:stieltjes}
\zeta(1+t,a)=\frac1t+\sum_{r\ge0}\frac{(-1)^r\gamma_r(a)}{r!}t^r,
\qquad \gamma_0(a)=-\psi(a).
\end{equation}
The generalized harmonic numbers are
\[
H_m^{(r)}=\sum_{k=1}^m k^{-r},\qquad H_m=H_m^{(1)},\qquad H_0^{(r)}=0.
\]
The complete exponential Bell polynomials $Y_r$ are defined by
\begin{equation}\label{hqj:eq:bell}
\exp\left(\sum_{j\ge1}\kappa_j\frac{t^j}{j!}\right)
=\sum_{r\ge0}Y_r(\kappa_1,\ldots,\kappa_r)\frac{t^r}{r!}.
\end{equation}
A superscript $[1]$ on $\Li_p$ means $\partial_p\Li_p$, whereas
$\psi^{(j)}$ denotes the ordinary $j$th derivative of digamma.

The notation $[x^{-1}]$ extracts the coefficient in the expansion at
positive infinity. A finite part is never an unspecified summation
method: the regularizing coordinate and the subtraction are displayed.
All new numerical comparisons are floating-point diagnostics. The
proofs, not those comparisons, establish the all-index statements.
